\chapter{Models and Methods}
\label{chap:02}
\setlength{\parskip}{1em}

\section{Research Design}
\label{sec:research_design}

\subsection{Problem Formulation}

This study is formulated as a \textbf{comparative supervised multi-label
classification experiment} on 12-lead electrocardiogram (ECG) signals.
The clinical motivation is straightforward: a single ECG recording may
simultaneously exhibit multiple pathological patterns.
For example, a patient with left ventricular hypertrophy (HYP) may also
present with conduction disturbances (CD) or ST/T-wave changes (STTC)
as secondary findings.
A binary-relevance (multi-label) formulation is therefore more faithful
to clinical reality than a mutually exclusive (multi-class) assignment,
which would force the model to choose a single dominant diagnosis
\cite{ref19,ref7}.

\subsection*{Task Definition}

Given a 10-second, 12-lead ECG signal sampled at 100\,Hz, the model
must predict a \emph{binary label vector}
$\mathbf{y} \in \{0,1\}^{5}$ whose components correspond to the five
PTB-XL diagnostic superclasses:

\begin{itemize}
  \item \textbf{NORM} --- Normal sinus rhythm
  \item \textbf{MI}   --- Myocardial Infarction
  \item \textbf{STTC} --- ST/T-Change
  \item \textbf{CD}   --- Conduction Disturbance
  \item \textbf{HYP}  --- Hypertrophy
\end{itemize}

\noindent Each label $y_c \in \{0,1\}$ is treated as an independent
binary classification sub-problem; a record may carry zero, one, or
several positive labels simultaneously.
This \emph{binary-relevance} decomposition transforms the joint
multi-label problem into $C = 5$ parallel binary tasks, each optimised
with its own sigmoid output and Binary Cross-Entropy objective.

\subsection*{Evaluation Protocol}

All five architectures are evaluated under an identical protocol
implemented in a shared \texttt{PTBXLEvaluator} module to ensure strict
comparability.
The primary metric is \emph{macro-averaged AUROC} (Area Under the
Receiver Operating Characteristic curve), which is the standard
benchmark metric for PTB-XL \cite{ref6}.
AUROC is threshold-free and measures the probability that the model
ranks a randomly chosen positive example above a randomly chosen
negative example; macro-averaging weights all five classes equally
regardless of their prevalence.

Supplementary metrics reported are:

\begin{itemize}
  \item \textbf{Micro AUROC} --- pooled over all labels and samples,
        giving larger classes more influence.
  \item \textbf{Macro AUPRC} --- macro-averaged Area Under the
        Precision-Recall Curve, informative for the class-imbalanced
        HYP and CD categories.
  \item $\mathbf{F_{\max}}$ --- maximum F1-score over all classification
        thresholds, found by sweeping $\tau \in [0,1]$ and reporting
        $\max_\tau F_1(\tau)$.
  \item \textbf{Sensitivity} (Recall) and \textbf{Specificity} at the
        $F_{\max}$ threshold.
\end{itemize}

\noindent A fixed random seed of 42 is used throughout all experiments
to guarantee reproducibility across data loading, model initialisation,
and augmentation.
The test fold is held out and evaluated only once, after all
hyperparameter decisions are finalised on the validation fold.

\section{Data Collection}
\label{sec:data_collection}

\subsection{PTB-XL Dataset}

All experiments are conducted on \textbf{PTB-XL} (Wagner et al., 2020),
the largest publicly available clinical 12-lead ECG dataset to date
\cite{ref12}.
PTB-XL contains 21{,}799 recordings collected from 18{,}869 patients at
Charit\'{e} -- Universit\"{a}tsmedizin Berlin between 1989 and 1996.
Each record is 10 seconds long and stored in WFDB format.
Annotations were performed by up to two cardiologists using standardised
SCP-ECG codes; in cases of disagreement, a validated automated algorithm
provided a third opinion \cite{ref12,ref16}.

The dataset is notable for three characteristics that make it the
de-facto standard for ECG deep learning benchmarking:
\begin{enumerate}
  \item \textbf{Scale}: at 21{,}799 recordings it is an order of
        magnitude larger than most prior ECG datasets.
  \item \textbf{Label quality}: multi-annotator SCP codes with
        likelihood scores allow probabilistic label assignment.
  \item \textbf{Official benchmark split}: a stratified ten-fold
        cross-validation protocol is distributed with the dataset,
        enabling reproducible comparison across research groups.
\end{enumerate}

Key dataset characteristics are summarised in Table~\ref{tab:ptbxl}.

\begin{table}[htbp]
  \centering
  \caption{PTB-XL dataset characteristics}
  \label{tab:ptbxl}
  \begin{tabular}{ll}
    \hline
    \textbf{Property}           & \textbf{Value}                    \\
    \hline
    Total recordings            & 21{,}799                          \\
    Patients                    & 18{,}869                          \\
    Leads                       & 12 (standard)                     \\
    Duration per record         & 10 seconds                        \\
    Sampling rate used          & 100\,Hz                           \\
    Sampling rate available     & 500\,Hz                           \\
    Label encoding              & SCP-codes $\to$ 5 superclasses    \\
    Annotation                  & 2 cardiologists + algorithm       \\
    Format                      & WFDB (PhysioNet)                  \\
    \hline
  \end{tabular}
\end{table}

\subsection{Label Encoding}

Raw SCP-codes are mapped to five diagnostic superclasses via the
\texttt{scp\_statements.csv} file distributed with the dataset.
Each SCP-code entry includes a field \texttt{diagnostic\_class} which
identifies the superclass, and a likelihood score in $[0,100]$.
A record is assigned to superclass $c$ if at least one of its SCP-codes
maps to $c$ with likelihood $> 0$.
Formally, let $\mathcal{S}(r)$ be the set of SCP-codes for record $r$
and $\ell(s)$ the likelihood score of code $s$:
\begin{equation}
  y_c^{(r)} = \mathbf{1}\!\left[
    \exists\, s \in \mathcal{S}(r)\colon
    \text{class}(s) = c \;\wedge\; \ell(s) > 0
  \right].
  \label{eq:label}
\end{equation}

The mapping is performed using \texttt{MultiLabelBinarizer}, producing a
binary vector $\mathbf{y} \in \{0,1\}^{5}$ per record.
Table~\ref{tab:class_dist} shows the resulting class distribution.

\begin{table}[htbp]
  \centering
  \caption{Class distribution in PTB-XL (diagnostic superclasses)}
  \label{tab:class_dist}
  \begin{tabular}{lrr}
    \hline
    \textbf{Class} & \textbf{Positive samples} & \textbf{Prevalence (\%)} \\
    \hline
    NORM  & 9{,}528  & 43.7 \\
    MI    & 5{,}486  & 25.2 \\
    STTC  & 5{,}250  & 24.1 \\
    CD    & 4{,}907  & 22.5 \\
    HYP   & 2{,}655  & 12.2 \\
    \hline
    \multicolumn{3}{l}{\footnotesize Totals exceed 100\% due to multi-label co-occurrence.}
  \end{tabular}
\end{table}

\noindent The severe class imbalance, particularly for HYP
(12.2\,\%), is a major challenge.
AUROC and AUPRC are chosen as primary metrics precisely because they are
robust to class imbalance; threshold-based metrics such as accuracy are
misleading under such skewed priors.

\subsection{Data Splits}

PTB-XL provides a stratified ten-fold split via the \texttt{strat\_fold}
field, where stratification ensures each fold contains a proportional
representation of all superclasses.
We follow the official benchmark protocol \cite{ref6}:

\begin{table}[htbp]
  \centering
  \caption{Train / validation / test split following the PTB-XL benchmark}
  \label{tab:splits}
  \begin{tabular}{lll}
    \hline
    \textbf{Subset} & \textbf{Folds} & \textbf{Purpose}            \\
    \hline
    Train           & 1--8           & Model optimisation          \\
    Validation      & 9              & Hyperparameter selection    \\
    Test            & 10             & Final evaluation (held-out) \\
    \hline
  \end{tabular}
\end{table}

\section{Data Preprocessing}
\label{sec:preprocessing}

The preprocessing pipeline is implemented in the
\texttt{PTBXLPreprocessorSOTA} class and applied identically to all three
subsets before any model-specific processing.
The pipeline consists of seven sequential steps described below and
illustrated schematically in Figure~\ref{fig:preprocess_pipeline}.

\begin{figure}[htbp]
  \centering
  \fbox{\parbox{0.85\linewidth}{\centering\vspace{2cm}
    \textit{[Место для рисунка: Pipeline предобработки сигнала ЭКГ]}
  \vspace{2cm}}}
  \caption{Preprocessing pipeline applied to each PTB-XL recording.
    Steps are applied in the order shown; the output is a cached
    \texttt{float32} array of shape $(12, 1000)$ per record.}
  \label{fig:preprocess_pipeline}
\end{figure}

\subsection{Signal Loading and Transposition}

Raw WFDB records are read using the \texttt{wfdb} library, yielding a
$1000 \times 12$ array (time $\times$ leads at 100\,Hz).
The array is transposed to $12 \times 1000$ (leads $\times$ time) to
conform to PyTorch's channel-first convention required by 1-D
convolutional architectures \cite{ref29,ref21}.
All subsequent operations act on the $12 \times 1000$ tensor.

\subsection{Length Normalisation}

Each recording is verified to contain exactly 1{,}000 samples
(10\,s at 100\,Hz).
Shorter sequences are padded to 1{,}000 samples using edge-padding
(\texttt{np.pad}, \texttt{mode=`edge'});
longer sequences are truncated to 1{,}000 samples from the start.
Edge-padding replicates the last valid sample, which introduces
less spectral distortion than zero-padding.

\subsection{Bandpass Filtering}

A 4th-order Butterworth bandpass filter with passband
$[0.5,\, 40]\,\text{Hz}$ is applied to each lead independently using
zero-phase forward--backward filtering (\texttt{scipy.signal.filtfilt}).
The transfer function of a 4th-order Butterworth filter is:
\begin{equation}
  |H(j\Omega)|^2 = \frac{1}{1 + \left(\Omega / \Omega_c\right)^{2N}},
  \quad N = 4,
  \label{eq:butterworth}
\end{equation}
where $\Omega_c$ is the $-3\,\text{dB}$ cutoff frequency.
The lower cutoff at 0.5\,Hz removes baseline wander caused by
respiratory motion and patient movement;
the upper cutoff at 40\,Hz suppresses high-frequency noise and
power-line interference while retaining clinically relevant waveform
morphology (QRS complexes, P waves, T waves lie within $[0.5, 40]\,\text{Hz}$)
\cite{ref37,ref23}.
Forward--backward (\texttt{filtfilt}) application doubles the effective
filter order to 8, achieving zero phase distortion at the cost of
increased computational load.

\subsection{Baseline Removal}

After bandpass filtering, the per-lead median is subtracted to eliminate
any residual DC offset:
\begin{equation}
  \hat{x}_{l,t} = x_{l,t} - \mathrm{median}_t(x_{l,\cdot}).
  \label{eq:baseline}
\end{equation}
Using the median rather than the mean makes this step robust to
impulsive noise (e.g.\ electrode pop artefacts), which would
artificially shift the mean.

\subsection{Z-Score Normalisation}

Each lead is independently standardised to zero mean and unit variance:
\begin{equation}
  \tilde{x}_{l,t} =
    \frac{x_{l,t} - \mu_l}{\sigma_l + \varepsilon},
  \qquad \varepsilon = 10^{-8},
  \label{eq:zscore}
\end{equation}
where
\begin{equation}
  \mu_l = \frac{1}{T}\sum_{t=1}^{T} x_{l,t},
  \qquad
  \sigma_l = \sqrt{\frac{1}{T}\sum_{t=1}^{T}(x_{l,t} - \mu_l)^2},
  \quad T = 1000.
  \label{eq:zscore_stats}
\end{equation}
The additive $\varepsilon$ prevents division by zero for flat or
near-flat leads.
Per-lead normalisation rather than global normalisation is used because
ECG leads have systematically different amplitude scales (e.g.\ limb
leads vs.\ precordial leads), and cross-lead amplitude ratios carry
diagnostic information that global normalisation would distort.

\subsection{Outlier Clipping}

Finally, values outside $[-10\sigma_l,\, +10\sigma_l]$ (evaluated before
normalisation) are clipped to suppress extreme artefacts without
discarding the recording entirely:
\begin{equation}
  x_{l,t}^{\mathrm{clipped}} = \mathrm{clip}(x_{l,t},\,
    \mu_l - 10\sigma_l,\, \mu_l + 10\sigma_l).
  \label{eq:clip}
\end{equation}
This step is applied \emph{before} Z-score normalisation in the
implementation, so the 10-sigma threshold is computed on the
baseline-corrected signal.

\subsection{Training-Time Augmentations}

Five stochastic augmentations are applied \emph{only to the training set}
during data loading.
Each augmentation is applied independently with probability $p = 0.5$,
following best practices for ECG deep learning \cite{ref7,ref16}.
The augmentations are composed using a pipeline:
\begin{equation}
  x_{\mathrm{aug}} = T_5 \circ T_4 \circ T_3 \circ T_2 \circ T_1 (x),
  \quad T_i \text{ applied with prob.\ } p = 0.5.
  \label{eq:augpipeline}
\end{equation}

\begin{table}[htbp]
  \centering
  \caption{Data augmentations applied during training}
  \label{tab:augmentations}
  \begin{tabular}{p{3.5cm}p{8.5cm}}
    \hline
    \textbf{Augmentation}  & \textbf{Description} \\
    \hline
    Gaussian Noise    &
      Additive white Gaussian noise:
      $\tilde{x} = x + \mathcal{N}(0,\, \sigma^2 \mathbf{I})$,
      $\sigma = 0.02$; simulates electrode measurement noise. \\[4pt]
    Amplitude Scaling &
      Per-lead multiplicative scale factor:
      $\tilde{x}_l = \alpha \cdot x_l$,
      $\alpha \sim \mathcal{U}(0.85,\, 1.15)$;
      simulates electrode contact variability. \\[4pt]
    Baseline Shift    &
      Per-lead additive offset:
      $\tilde{x}_l = x_l + \delta_l$,
      $\delta_l \sim \mathcal{U}(-0.05,\, +0.05)$;
      simulates residual baseline wander. \\[4pt]
    Time Shift        &
      Circular shift by a random integer $s \sim \mathcal{U}[-30, +30]$
      samples:
      $\tilde{x}_l[t] = x_l[(t - s)\!\mod T]$;
      simulates R-peak trigger jitter. \\[4pt]
    Lead Dropout      &
      Each lead is independently zeroed with probability $p = 0.05$:
      $\tilde{x}_l = m_l \cdot x_l$,
      $m_l \sim \mathrm{Bernoulli}(0.95)$;
      encourages robustness to missing or noisy channels \cite{ref22}. \\
    \hline
  \end{tabular}
\end{table}

\noindent Crucially, augmentations are applied on-the-fly during
DataLoader iteration, so each training epoch sees a different
stochastic realisation of every sample.
Validation and test sets receive no augmentation;
Test-Time Augmentation (TTA) used for the RetNet model is described
separately in Section~\ref{subsec:model_d}.

\section{Model Development}
\label{sec:model_development}

Five neural network architectures of increasing complexity and
inductive bias were developed and evaluated.
All share the same training configuration (Table~\ref{tab:training})
and differ only in their architectural components and inductive biases.
The five models span a spectrum from a plain CNN baseline
(Model~A, \texttildelow 0.4M parameters) to a full Retention-based
sequence model with cross-lead fusion
(Model~D / RetNet-ECG, 33.5M parameters).

\begin{table}[htbp]
  \centering
  \caption{Common training hyperparameters for all models}
  \label{tab:training}
  \begin{tabular}{ll}
    \hline
    \textbf{Hyperparameter} & \textbf{Value}                           \\
    \hline
    Optimiser               & Adam ($\beta_1{=}0.9$, $\beta_2{=}0.999$) \\
    Learning rate schedule  & Linear warmup (5 epochs) + cosine decay  \\
    Peak learning rate      & $10^{-3}$                                \\
    Batch size              & 64                                       \\
    Maximum epochs          & 50                                       \\
    Weight decay            & $10^{-4}$                                \\
    Loss function           & Binary Cross-Entropy (BCE)               \\
    Early stopping          & Patience 10 (macro AUROC on val)         \\
    Random seed             & 42                                       \\
    \hline
  \end{tabular}
\end{table}

\noindent \textbf{Loss function.}
Binary Cross-Entropy is the standard loss for independent multi-label
problems.
For a batch of $B$ samples and $C = 5$ classes it decomposes as:
\begin{equation}
  \mathcal{L}_{\mathrm{BCE}} =
    -\frac{1}{BC}\sum_{i=1}^{B}\sum_{c=1}^{C}
    \Bigl[
      y_{ic}\log\sigma(\hat{z}_{ic})
      + (1 - y_{ic})\log\bigl(1 - \sigma(\hat{z}_{ic})\bigr)
    \Bigr],
  \label{eq:bce}
\end{equation}
where $\hat{z}_{ic}$ is the raw logit for sample $i$ and class $c$,
and $\sigma(\cdot)$ is the sigmoid function \cite{ref19,ref24}.

\noindent \textbf{Learning rate schedule.}
The learning rate follows a linear warmup for the first 5 epochs from
$10^{-6}$ to $10^{-3}$, followed by cosine annealing to $10^{-6}$:
\begin{equation}
  \eta_t =
    \eta_{\min}
    + \tfrac{1}{2}(\eta_{\max} - \eta_{\min})
    \left(1 + \cos\!\frac{\pi(t - T_{\mathrm{warm}})}{T_{\max} - T_{\mathrm{warm}}}\right),
  \label{eq:cosine}
\end{equation}
for $t > T_{\mathrm{warm}}$, where $T_{\max} = 50$ epochs and
$\eta_{\max} = 10^{-3}$.

\subsection{Baseline --- ResNet-1D}
\label{subsec:model_a}

\subsubsection{Architecture}

The model accepts input tensors of shape $(B, 12, 1000)$ (batch,
leads, time samples).
A \emph{reshape trick} is used: the 12-lead signal is reshaped to
$(B \times 12, 1, 1000)$, treating each lead as an independent
1-channel time series.
This allows the backbone to use shared weights across all leads,
effectively $12\times$ data augmentation at no parameter cost.

The backbone consists of five \texttt{ResidualBlock1d} blocks with the
following channel progression:
\begin{equation}
  \underbrace{1}_{\text{input}} \to
  \underbrace{32}_{\text{stem}} \to
  \underbrace{64}_{\text{B1}} \to
  \underbrace{128}_{\text{B2}} \to
  \underbrace{256}_{\text{B3}} \to
  \underbrace{256}_{\text{B4}},
  \label{eq:cnn_channels}
\end{equation}
with strides $(1, 2, 2, 2, 2)$ so the temporal dimension is reduced
from 1000 to approximately 63.
Each residual block follows the structure:
\begin{equation}
  \mathbf{h} = \mathrm{ReLU}\!\left(
    \mathrm{SE}\!\left(
      \mathrm{BN}(\mathrm{Conv}_7(\mathbf{x}))
    \right) + W_s\mathbf{x}
  \right),
  \label{eq:resblock}
\end{equation}
where $\mathrm{Conv}_7$ denotes a 1-D convolution with kernel size 7
(receptive field $\approx 70\,\text{ms}$ at 100\,Hz, sufficient to
capture a complete QRS complex), $\mathrm{BN}$ is Batch Normalisation,
$\mathrm{SE}$ is the Squeeze-and-Excitation recalibration (see
Section~\ref{subsec:model_b}), and $W_s$ is a $1\times 1$ projection
to match dimensions when the channel count or stride changes.

After the backbone, Global Average Pooling (GAP) collapses the temporal
dimension:
\begin{equation}
  \mathbf{g}_l = \frac{1}{T'}\sum_{t=1}^{T'} \mathbf{h}_{l,t}
  \in \mathbb{R}^{256},
  \label{eq:gap}
\end{equation}
where $T'$ is the temporal length after striding.
The per-lead embeddings are averaged across all 12 leads to obtain a
single 256-dimensional vector, which is passed through a two-layer MLP
classifier:
\begin{equation}
  \hat{\mathbf{z}} = W_2\, \mathrm{GELU}(W_1 \mathbf{g} + \mathbf{b}_1) + \mathbf{b}_2
  \in \mathbb{R}^5.
  \label{eq:mlp}
\end{equation}

\subsubsection{Architecture Figure}

\begin{figure}[htbp]
  \centering
  \fbox{\parbox{0.85\linewidth}{\centering\vspace{2cm}
    \textit{[Место для рисунка: Архитектура Model A --- CNN Baseline]}
  \vspace{2cm}}}
  \caption{Model A (CNN Baseline) architecture.
    Input $(B, 12, 1000)$ is reshaped to $(B{\times}12, 1, 1000)$
    and processed by five ResidualBlock1d stages with increasing
    channel counts. Global average pooling produces per-lead embeddings;
    lead-mean pooling and a two-layer MLP yield five logits.}
  \label{fig:model_a}
\end{figure}

\subsection{Baseline --- InceptionTime}
\label{subsec:inception}

\subsubsection{Architecture}

Each \textbf{Inception module} applies three parallel 1-D convolutions
with kernel sizes $k \in \{10, 20, 40\}$ samples alongside a
MaxPool--Conv shortcut branch.
A $1\times 1$ bottleneck convolution first reduces the channel count to
32 before the parallel branches, limiting parameter growth:
\begin{equation}
  \mathbf{h} = \mathrm{BN}\!\left(
    \mathrm{Cat}\!\left[
      \mathrm{Conv}_{10}(\mathbf{b}),\;
      \mathrm{Conv}_{20}(\mathbf{b}),\;
      \mathrm{Conv}_{40}(\mathbf{b}),\;
      \mathrm{Conv}_1(\mathrm{MaxPool}_3(\mathbf{x}))
    \right]
  \right),
  \label{eq:inception_module}
\end{equation}
where $\mathbf{b} = \mathrm{Conv}_1(\mathbf{x})$ is the bottleneck
output.
All convolutions use \texttt{padding=`same'} to preserve temporal length.
The four branches each produce $n = 32$ filters, so the module outputs
$4n = 128$ channels.

\textbf{Inception blocks} stack three consecutive Inception modules with
a residual connection:
\begin{equation}
  \mathbf{y}_k = \mathrm{ReLU}\!\left(
    F_k(\mathbf{y}_{k-1}) + W_s\,\mathbf{y}_{k-1}
  \right),
  \label{eq:inception_block}
\end{equation}
where $F_k$ is the $k$-th depth-3 stack of Inception modules and
$W_s$ is a shortcut $1\times 1$ convolution.
Two such Inception blocks are stacked in total.

The same reshape trick as Model~A is applied: input $(B, 12, 1000)$ is
reshaped to $(B\times 12, 1, 1000)$, each lead processed independently.
After the encoder, Global Average Pooling followed by lead-mean pooling
produces a single feature vector for the linear classifier.

\subsubsection{Architecture Figure}

\begin{figure}[htbp]
  \centering
  \fbox{\parbox{0.85\linewidth}{\centering\vspace{2cm}
    \textit{[Место для рисунка: Архитектура InceptionTime]}
  \vspace{2cm}}}
  \caption{InceptionTime architecture.
    Four parallel branches per Inception module capture morphological
    patterns at scales 10, 20, and 40 samples, plus a MaxPool shortcut.
    Two Inception blocks with residual connections are stacked before
    global average pooling and a linear classifier.}
  \label{fig:inception}
\end{figure}

\subsection{Model B --- Squeeze-and-Excitation Channel Attention}
\label{subsec:model_b}

\subsubsection{Squeeze-and-Excitation Block}

The SE block \cite{ref5} is inserted after each convolutional stage.
It operates in two steps:

\paragraph{Squeeze.}
Global average pooling compresses each feature map channel $c$ of width
$T'$ to a scalar descriptor:
\begin{equation}
  s_c = \frac{1}{T'}\sum_{t=1}^{T'} h_{c,t}.
  \label{eq:squeeze}
\end{equation}

\paragraph{Excitation.}
The channel descriptor vector $\mathbf{s} \in \mathbb{R}^C$ is passed
through a two-layer bottleneck MLP with reduction ratio $r$:
\begin{equation}
  \mathbf{e} = \sigma\!\left(
    W_2\, \mathrm{ReLU}(W_1 \mathbf{s})
  \right), \quad
  W_1 \in \mathbb{R}^{C/r \times C},\;
  W_2 \in \mathbb{R}^{C \times C/r},\;
  r = 8.
  \label{eq:excitation}
\end{equation}

\paragraph{Recalibration.}
Each channel of the feature map is scaled by its excitation scalar:
\begin{equation}
  \tilde{h}_{c,t} = e_c \cdot h_{c,t}.
  \label{eq:recalibration}
\end{equation}
The SE block adds only $2C^2/r$ parameters per stage.
With $C=256$ and $r=8$ this is $\approx 16{,}384$ parameters---less than
1\,\% of the backbone---while yielding a measurable AUROC improvement
\cite{ref8,ref36}.

\subsubsection{Lead Attention Pooling}

After the CNN+SE backbone produces per-lead embeddings
$\{\mathbf{g}_l\}_{l=1}^{12}$, the 12 embeddings are pooled with
learned attention weights instead of simple averaging.

A global context vector is first computed:
\begin{equation}
  \mathbf{c} = \frac{1}{12}\sum_{l=1}^{12}\mathbf{g}_l \in \mathbb{R}^{256}.
  \label{eq:global_ctx}
\end{equation}
Each lead embedding is concatenated with this global context and passed
through a small MLP to produce a scalar attention score:
\begin{equation}
  a_l = f_\mathrm{att}\!\left(
    \mathrm{LayerNorm}([\mathbf{g}_l;\, \mathbf{c}])
  \right),\quad
  f_\mathrm{att}: \mathbb{R}^{512} \to \mathbb{R}.
  \label{eq:lead_att_score}
\end{equation}
The scores are normalised by a learnable temperature $\tau$:
\begin{equation}
  \alpha_l = \frac{\exp(a_l / \tau)}{\sum_{l'=1}^{12}\exp(a_{l'} / \tau)}.
  \label{eq:lead_softmax}
\end{equation}
The attention-weighted embedding, together with max-pooled and
mean-pooled embeddings, forms a \emph{generalised pooling} vector:
\begin{equation}
  \mathbf{p} = \left[
    \sum_{l=1}^{12}\alpha_l \mathbf{g}_l
    \;;\;
    \max_l \mathbf{g}_l
    \;;\;
    \frac{1}{12}\sum_{l=1}^{12}\mathbf{g}_l
  \right] \in \mathbb{R}^{768}.
  \label{eq:gen_pool}
\end{equation}
This concatenation is fed to the two-layer MLP classifier.

\subsubsection{Architecture Figure}

\begin{figure}[htbp]
  \centering
  \fbox{\parbox{0.85\linewidth}{\centering\vspace{2cm}
    \textit{[Место для рисунка: Архитектура Model B --- SE Attention]}
  \vspace{2cm}}}
  \caption{Model B (SE Channel Attention) architecture.
    The CNN encoder is identical to Model A but with SE blocks after
    each residual stage. Lead-wise embeddings are pooled via
    context-aware soft attention; the attention-pooled, max-pooled,
    and mean-pooled embeddings are concatenated before classification.}
  \label{fig:model_b}
\end{figure}

\subsection{Model C --- LeadWise Graph Neural Network}
\label{subsec:model_c}

\subsubsection{Stage 1: Lead Encoding}

Each of the 12 leads is independently encoded by a shared
\texttt{LeadWiseResNet1d} backbone---a lead-adapted version of
ResNet1dWang \cite{ref34}---producing a $d = 128$-dimensional embedding:
\begin{equation}
  \mathbf{H} = \mathrm{Encoder}(X) \in \mathbb{R}^{12 \times 128},
  \label{eq:lead_embed}
\end{equation}
where row $l$ of $\mathbf{H}$ is the embedding of lead $l$.
The encoder uses the same shared weights for all 12 leads
(parameter sharing), ensuring translational equivariance across leads.

\subsubsection{Clinical Adjacency Graph}

An adjacency matrix $\mathbf{A} \in \{0,1\}^{12 \times 12}$ encodes
known electrophysiological relationships among the 12 leads:

\begin{itemize}
  \item \textbf{Frontal leads} (I, II, III, aVR, aVL, aVF): fully
        connected, reflecting Kirchhoff's voltage law
        ($\mathrm{I} = \mathrm{II} - \mathrm{III}$,
        $\mathrm{aVR} = -(I + II)/2$).
  \item \textbf{Precordial leads} (V1--V6): connected in a chain of
        neighbours plus one-skip connections
        ($V_i \leftrightarrow V_{i+2}$), reflecting their
        spatial arrangement around the chest.
  \item \textbf{Cross-group edges}: clinically meaningful pairs
        (I\,↔\,V5, I\,↔\,V6, II\,↔\,V4, aVF\,↔\,V4,
        aVL\,↔\,V5, aVR\,↔\,V1).
\end{itemize}

The raw adjacency matrix is symmetrically normalised following
Kipf \& Welling \cite{ref20}:
\begin{equation}
  \hat{\mathbf{A}} = \tilde{\mathbf{D}}^{-1/2}
    (\mathbf{A} + \mathbf{I})
    \tilde{\mathbf{D}}^{-1/2},
  \label{eq:gcn_norm}
\end{equation}
where $\tilde{\mathbf{D}}_{ii} = \sum_j (A_{ij} + \delta_{ij})$ is the
degree matrix of the adjacency with self-loops added.

\subsubsection{Graph Attention Network}

Two layers of Graph Attention Network (GAT) \cite{ref24} propagate
information across the graph.
For a single GAT layer, the attention coefficient between nodes $i$ and
$j$ is:
\begin{equation}
  e_{ij} = \mathrm{LeakyReLU}\!\left(
    \mathbf{a}^\top
    \left[ W\mathbf{h}_i \;\|\; W\mathbf{h}_j \right]
  \right),
  \label{eq:gat_energy}
\end{equation}
where $W \in \mathbb{R}^{d_h \times d}$ is a learnable weight matrix,
$\mathbf{a} \in \mathbb{R}^{2d_h}$ is an attention vector, and
$\|$ denotes concatenation.
The coefficient is normalised over the neighbourhood
$\mathcal{N}(i) = \{j : A_{ij} > 0\} \cup \{i\}$:
\begin{equation}
  \alpha_{ij} = \frac{\exp(e_{ij})}{\sum_{k \in \mathcal{N}(i)}\exp(e_{ik})}.
  \label{eq:gat_softmax}
\end{equation}
The updated node embedding is the attention-weighted sum over neighbours,
with multi-head aggregation (4 heads in layer 1, 1 head in layer 2):
\begin{equation}
  \mathbf{h}_i' = \Big\|_{m=1}^{M}
    \mathrm{ELU}\!\left(
      \sum_{j \in \mathcal{N}(i)} \alpha_{ij}^{(m)} W^{(m)} \mathbf{h}_j
    \right),
  \label{eq:gat_update}
\end{equation}
where $M$ is the number of attention heads.
Residual connections are added after each GAT layer:
\begin{equation}
  \mathbf{h}^{(k+1)} = \mathrm{LayerNorm}\!\left(
    \mathbf{h}^{(k+1)}_\mathrm{GAT} + W_r\, \mathbf{h}^{(k)}
  \right).
  \label{eq:gat_residual}
\end{equation}

\subsubsection{Per-Class Attention Readout}

A key feature of the architecture is a per-class attention readout that
allows each diagnostic class $c$ to attend selectively to the leads most
informative for that class.
Learnable query vectors $\mathbf{q}_c \in \mathbb{R}^{d_h}$ are used:
\begin{equation}
  s_{c,l} = \frac{\mathbf{q}_c^\top \mathbf{h}_l^{(2)}}{\sqrt{d_h}},
  \quad
  \beta_{c,l} = \frac{\exp(s_{c,l})}{\sum_{l'}\exp(s_{c,l'})},
  \label{eq:class_readout}
\end{equation}
\begin{equation}
  \mathbf{p}_c = \sum_{l=1}^{12} \beta_{c,l}\, \mathbf{h}_l^{(2)}.
  \label{eq:class_pool}
\end{equation}
A two-layer MLP then maps $\mathbf{p}_c$ to a scalar logit $\hat{z}_c$.
The attention weights $\{\beta_{c,l}\}$ provide explicit per-class
lead-importance maps that can be used for interpretability
(explainable AI output).

\subsubsection{Two-Stage Training}

To stabilise GNN training, the model is trained in two stages:
\begin{enumerate}
  \item \textbf{Stage 1}: Encoder weights are initialised from the
        pre-trained ResNet1dWang baseline; the adjacency matrix
        $\mathbf{A}$ is held fixed.
        Only the GNN head is updated.
  \item \textbf{Stage 2}: All components (encoder + GNN head + adjacency)
        are jointly fine-tuned.
\end{enumerate}
Predicted probabilities are subsequently calibrated using
\emph{temperature scaling}: a single scalar $T_\mathrm{cal}$ is learned
on the validation set such that $\sigma(\hat{z}_c / T_\mathrm{cal})$
is better calibrated than $\sigma(\hat{z}_c)$ \cite{ref24}.

\subsubsection{Architecture Figure}

\begin{figure}[htbp]
  \centering
  \fbox{\parbox{0.85\linewidth}{\centering\vspace{2.5cm}
    \textit{[Место для рисунка: Архитектура Model C --- LeadWise GNN
    (Dimash пришлёт схему)]}
  \vspace{2.5cm}}}
  \caption{Model C (LeadWise GNN) architecture.
    A shared LeadWiseResNet1d encoder produces per-lead embeddings.
    Two GAT layers with residual connections propagate information
    along the clinical adjacency graph.
    A per-class attention readout produces class-specific pooled
    representations before the MLP classifier.}
  \label{fig:model_c}
\end{figure}

\subsection{Model D --- RetNet-ECG}
\label{subsec:model_d}

\subsubsection{MultiScalePatchEmbed}

The raw 12-lead input $(B, 12, 1000)$ is projected into a token
sequence via an Inception-style patch embedding layer with three
parallel convolutional kernels of sizes $k \in \{5, 11, 25\}$ samples
and stride $s = 5$:
\begin{equation}
  \mathbf{E} = \mathrm{Conv}_1\!\left(
    \mathrm{Cat}\!\left[
      \mathrm{GELU}(\mathrm{Conv}_{5}(x)),\;
      \mathrm{GELU}(\mathrm{Conv}_{11}(x)),\;
      \mathrm{GELU}(\mathrm{Conv}_{25}(x))
    \right]
  \right) \in \mathbb{R}^{T \times d},
  \label{eq:patch_embed}
\end{equation}
where $T = 200$ tokens per lead ($(1000 - 5)/5 + 1 = 200$),
$d = d_{\mathrm{model}} = 512$.
The three branches encode:
\begin{itemize}
  \item $k{=}5$: fine-grained QRS morphology details
  \item $k{=}11$: QRS envelope structure
  \item $k{=}25$: P/T waves and slow oscillations
\end{itemize}
The outputs are summed channel-wise, yielding a sequence of
$D$-dimensional tokens that encode multiple temporal scales
simultaneously \cite{ref29,ref31}.

\subsubsection{xPos Positional Encoding}

Position-dependent scaling is injected through xPos rotary encodings
\cite{ref18}.
For query/key vectors in each retention head, the position-$t$
transformation is:
\begin{equation}
  \mathrm{xPos}(\mathbf{x}, t) =
  \begin{bmatrix}
    x_1 \cos(\theta_1 t) - x_2 \sin(\theta_1 t) \\
    x_1 \sin(\theta_1 t) + x_2 \cos(\theta_1 t) \\
    \vdots
  \end{bmatrix},
  \quad
  \theta_k = \frac{1}{10000^{2k/d_h}},
  \label{eq:xpos}
\end{equation}
where $d_h = d / H$ is the head dimension and $H$ is the number of
heads.

\subsubsection{MultiScaleRetention}

The Retention mechanism replaces self-attention with a linear-complexity
recurrent formulation that admits both parallel (training) and recurrent
(inference) computation modes \cite{ref18}.
For head $m$ with decay rate $\gamma_m$, the causal decay mask is:
\begin{equation}
  D_{ij}^{(m)} = \gamma_m^{i-j} \cdot \mathbf{1}[i \ge j],
  \label{eq:decay_mask}
\end{equation}
where $\gamma_m$ is defined by:
\begin{equation}
  \gamma_m = 1 - 2^{-(5 + m \cdot 3/H)},
  \quad m = 0, \ldots, H-1.
  \label{eq:gamma}
\end{equation}
This assigns different decay rates to different heads, allowing the
multi-head model to attend to both local ($\gamma \approx 0.5$) and
long-range ($\gamma \approx 0.97$) temporal contexts simultaneously.

The retention output for head $m$ in parallel mode is:
\begin{equation}
  \mathrm{Ret}_m(\mathbf{X}) =
  \frac{(\mathbf{Q}_m \mathbf{K}_m^\top \odot \mathbf{D}^{(m)})
        \cdot \mathbf{V}_m}
       {\max\!\left(
         \bigl|(\mathbf{D}^{(m)} \mathbf{1})\bigr|,\, 1
       \right)},
  \label{eq:retention}
\end{equation}
where $\mathbf{Q}_m = \mathbf{X} W_Q^{(m)}$,
$\mathbf{K}_m = \mathbf{X} W_K^{(m)}$,
$\mathbf{V}_m = \mathbf{X} W_V^{(m)}$,
all with xPos applied to $\mathbf{Q}_m$ and $\mathbf{K}_m$.
The denominator normalises by the sum of decay weights to prevent
output magnitude growth with sequence length.

Multi-scale retention combines $H$ heads with Group Normalisation and
a gated output projection:
\begin{equation}
  \mathrm{MSR}(\mathbf{X}) =
  W_O\!\left(
    \mathrm{GN}\!\left(
      \mathrm{Cat}_{m=1}^{H}\mathrm{Ret}_m(\mathbf{X})
    \right) \odot \mathrm{SiLU}(W_G \mathbf{X})
  \right),
  \label{eq:msr}
\end{equation}
where $\mathrm{GN}$ is Group Normalisation with $H$ groups and
$\mathrm{SiLU}(x) = x \cdot \sigma(x)$ is the Sigmoid Linear Unit.

\subsubsection{RetNetFFN with SwiGLU}

Each Retention block is followed by a Feed-Forward Network using the
SwiGLU activation \cite{ref18}:
\begin{equation}
  \mathrm{SwiGLU}(\mathbf{x}) =
    \mathrm{SiLU}(W_1 \mathbf{x}) \odot (W_2 \mathbf{x}),
  \label{eq:swiglu}
\end{equation}
followed by a linear projection $W_3$:
\begin{equation}
  \mathrm{FFN}(\mathbf{x}) = W_3\!\left(
    \mathrm{SiLU}(W_1 \mathbf{x}) \odot W_2 \mathbf{x}
  \right).
  \label{eq:ffn}
\end{equation}
The hidden dimension is $d_{\mathrm{ff}} = 2d_{\mathrm{model}} = 1024$.

\subsubsection{RetNetBlock with Stochastic Depth}

A RetNet block applies MSR and FFN with pre-normalisation and
DropPath (Stochastic Depth) regularisation:
\begin{align}
  \mathbf{x}' &= \mathbf{x} + \mathrm{DropPath}\!\left(
    \mathrm{MSR}(\mathrm{LN}(\mathbf{x}))
  \right), \label{eq:retblock1} \\
  \mathbf{x}'' &= \mathbf{x}' + \mathrm{DropPath}\!\left(
    \mathrm{FFN}(\mathrm{LN}(\mathbf{x}'))
  \right). \label{eq:retblock2}
\end{align}
DropPath stochastically drops entire residual paths during training:
\begin{equation}
  \mathrm{DropPath}(\mathbf{x}; p) =
  \begin{cases}
    \mathbf{x} / (1-p) & \text{with probability } 1-p, \\
    \mathbf{0}          & \text{with probability } p.
  \end{cases}
  \label{eq:droppath}
\end{equation}
Ten RetNet blocks are stacked with linearly increasing DropPath rates:
\begin{equation}
  p_k = p_{\max} \cdot \frac{k}{\max(N-1,\,1)},
  \quad k = 0, \ldots, N-1,\quad p_{\max} = 0.2.
  \label{eq:droppath_schedule}
\end{equation}

\subsubsection{LeadFusionTransformer}

After the temporal RetNet stack processes each lead independently,
temporal mean-pooling produces a lead-level representation:
\begin{equation}
  \mathbf{z}_l = \frac{1}{T}\sum_{t=1}^{T}\mathbf{x}_l^{(N)}
  \in \mathbb{R}^d, \quad l = 1,\ldots,12.
  \label{eq:temporal_pool}
\end{equation}
A two-layer cross-lead Transformer encoder aggregates information across
all 12 leads, treating leads as tokens:
\begin{equation}
  \mathbf{Z}' = \mathrm{TransformerEncoder}(\mathbf{Z})
  \in \mathbb{R}^{12 \times d},
  \label{eq:lead_fusion}
\end{equation}
where standard multi-head self-attention ($H = 8$ heads) is used.
This layer captures spatial inter-lead correlations---e.g.\ reciprocal
ST changes between inferior and lateral leads---that are beyond the
reach of single-lead processing \cite{ref24,ref31}.
A final lead-mean pooling yields the 512-dimensional classification
vector.

\subsubsection{Inference Enhancements}

Two complementary strategies are employed at inference time:

\begin{itemize}
  \item \textbf{ModelEMA} (Exponential Moving Average of weights,
        decay $= 0.9999$): Let $\theta_t$ be model parameters at
        training step $t$; the EMA maintains:
        \begin{equation}
          \tilde{\theta}_t = \delta \cdot \tilde{\theta}_{t-1}
            + (1 - \delta) \cdot \theta_t,
          \quad \delta = 0.9999.
          \label{eq:ema}
        \end{equation}
        At inference $\tilde{\theta}$ is used.
        EMA weights represent a temporally smoothed model, reducing
        prediction variance and yielding consistent AUROC improvements
        of +0.3--0.8\,\% \cite{ref10}.

  \item \textbf{Test-Time Augmentation (TTA$\times$15)}: Fifteen
        independently augmented versions of each test signal (using
        the same augmentation pipeline as training) are processed and
        their predicted probabilities are averaged:
        \begin{equation}
          \hat{p}_c = \frac{1}{15}\sum_{k=1}^{15}
            \sigma\!\left(\hat{z}_c^{(k)}\right).
          \label{eq:tta}
        \end{equation}
        This ensemble-within-a-model strategy yields consistent gains
        in AUROC without requiring multiple trained models
        \cite{ref7,ref16}.
\end{itemize}

\subsubsection{Architecture Figure}

\begin{figure}[htbp]
  \centering
  \includegraphics[width=\linewidth]{AITU_Diploma_2025-2026/chapters/chapter02/fig02/fig_2_1.png}
  \caption{Overview of the RetNet-ECG architecture. From left to
    right: \textit{MultiScalePatchEmbed} encodes the 12-lead input at
    three temporal scales ($k \in \{5,11,25\}$, stride 5, yielding
    $T=200$ tokens); ten stacked RetNet blocks with SwiGLU FFN and
    linearly increasing DropPath process each lead sequence;
    \textit{LeadFusionTransformer} (2-layer) aggregates cross-lead
    context via multi-head self-attention; global lead pooling and
    a two-layer MLP head produce five-class logits.
    At inference: ModelEMA weights (decay 0.9999) and
    TTA$\times$15 are applied.}
  \label{fig:retnet_architecture}
\end{figure}

\subsection{Summary of Architectures}
\label{subsec:arch_summary}

Table~\ref{tab:model_summary} provides a side-by-side overview of the
five architectures, ordered from simplest to most complex.

\begin{table}[htbp]
  \centering
  \caption{Summary of the five architectures compared in this study}
  \label{tab:model_summary}
  \begin{tabular}{p{2.8cm}p{3.2cm}p{3.8cm}r}
    \hline
    \textbf{Model}   & \textbf{Key Mechanism}
                     & \textbf{Distinguishing Feature}
                     & \textbf{Params}                              \\
    \hline
    CNN Baseline     & 1-D Conv + GAP
                     & Lower-bound reference
                     & $\sim$0.4M                                   \\
    InceptionTime    & Multi-kernel inception
                     & Multi-scale temporal features
                     & $\sim$0.5M                                   \\
    SE Attention     & SE block + lead attention
                     & Adaptive channel \& lead re-weighting
                     & $\sim$0.4M                                   \\
    LeadWise GNN     & GAT + clinical graph
                     & Learnable inter-lead topology
                     & $\sim$0.62M                                  \\
    RetNet-ECG    & Multi-scale Retention +
                       cross-lead Transformer
                     & Linear-complexity sequence model;
                       TTA$\times$15 \& EMA
                     & $\sim$33.5M                                  \\
    \hline
  \end{tabular}
\end{table}

